\section{Round-thirteen positive closure: the joint collision Gaussian field and its closed inverse form}
\label{sec:r13-b3}

The primitive noise lives on collision space.  Density and contact
fluctuations are two images of the same Gaussian random measure, and a dual
pair $(p,\psi)$ sees the single coefficient $\Delta p+\psi$.  Thus pure
contact tests retain their Poisson diagonal variance.

\subsection{The perturbative phase chart and energy space}

Fix a smooth Maxwellian $M$ and a compact chart of exposed paths satisfying
\[
 cM\le f(t,x,v)\le CM,
 \qquad
 \|f-M\|_{W^{2,\infty}(M^{-1})}
 +\|q-1\|_{W^{1,\infty}_w}\le\delta,
 \qquad 0<q_-\le q\le q_+.
\]
The chart is obtained by restricting the B2 real-source ball so that its first
and second pressure derivatives remain in these bounds.

Let $\mathsf V\subset\mathsf H\subset\mathsf V^*$ be the kinetic Gelfand
triple whose norm contains transport derivatives, the weighted collision
Dirichlet form, the five spatially dependent macroscopic coefficient fields,
and the B2 normal trace graph.

\begin{theorem}[Nonautonomous perturbative kinetic energy estimate]
\label{thm:r13-b3-energy}
For $\delta$ and $T$ sufficiently small, the linearized biased equation
\[
 \partial_tu+v\cdot\nabla_xu=\mathcal L_{f,q}u+F,
 \qquad u(0)=u_0,
\]
has a unique solution in
$L^2(0,T;\mathsf V)\cap C([0,T];\mathsf H)$ and
\[
 \sup_{t\le T}\|u_t\|_{\mathsf H}^2
 +c\int_0^T\|u_t\|_{\mathsf V}^2dt
 \le C\left(\|u_0\|_{\mathsf H}^2
 +\int_0^T\|F_t\|_{\mathsf V^*}^2dt\right).
\]
The trace pairing is continuous on this solution space.
\end{theorem}

\begin{proof}
For $q=1$ and $f=M$, the symmetric linearized collision form is coercive
modulo the five velocity invariants.  Projecting onto those invariants gives
spatially dependent mass, momentum, and energy fields.  Add the standard
transport--moment cross terms to the energy; integration by parts and the
moment equations control the macroscopic component without pretending it is
five scalar constants.  The $W^{2,\infty}$ coefficient perturbations and
$q-1$ terms are form-bounded with relative bound $O(\delta+T)$.  Absorb them
and use Lions' theorem for nonautonomous forms.  The B2 Green graph estimate
controls the normal trace by the graph part of $\mathsf V$.
\end{proof}

\subsection{The joint Gaussian random measure}

Let $\mathfrak C_T$ be collision space and define the intensity
\[
 d\mathfrak a=q(t,x,v,v_*,\omega)\,dA_f.
\]
Let $W$ be the centered isonormal Gaussian random measure on
$L^2(\mathfrak C_T,\mathfrak a)$, independent of the prepared initial
Gaussian field $G_0$.  For a density test $p$ and a direct contact test
$\psi$, write
\[
 W(p,\psi)=\int_{\mathfrak C_T}(\Delta p+\psi)\,dW.
\]

\begin{theorem}[Prepared joint density--contact Gaussian process]
\label{thm:r13-b3-joint}
There is a unique centered Gaussian pair $(u,\Xi)$ satisfying
\[
 d\langle u_t,p_t\rangle
 =\langle u_t,(\partial_t+v\cdot\nabla_x
              +\mathcal L_{f,q}^*)p_t\rangle dt
 +\int\Delta p_t\,dW_t,
\]
with initial law $G_0$, and
\[
 \Xi(dz)=W(dz)+D(qA_f)[u](dz).
\]
For two dual pairs $(p,\psi)$ and $(p',\psi')$,
\[
 d\langle M(p,\psi),M(p',\psi')\rangle_t
 =\int q(\Delta p+\psi)(\Delta p'+\psi')\,dA_{f_t}\,dt.
\]
In particular a pure contact test has variance
$\int q\psi^2dA_f>0$ unless it vanishes in $L^2(\mathfrak a)$.
\end{theorem}

\begin{proof}
The stochastic convolution with the evolution family from
Theorem~\ref{thm:r13-b3-energy} defines $u$.  Linearizing the collision
intensity $qA_f$ gives the drift response in $\Xi$, while the same random
measure $W$ is the diagonal fluctuation of actual contacts.  Applying a dual
pair combines the density jump and direct contact mark into
$\Delta p+\psi$.  The isometry of $W$ gives the bracket and uniqueness.
\end{proof}

\subsection{Microscopic cumulants and stopping-time tightness}

Let
\[
 u^\varepsilon=
 \sqrt{\mu_\varepsilon}(\pi^\varepsilon-f),
 \qquad
 \Xi^\varepsilon=
 \sqrt{\mu_\varepsilon}(\Gamma^\varepsilon-qA_f).
\]

\begin{lemma}[Conditional localized cumulants]
\label{lem:r13-b3-cumulants}
For every finite family of density/contact tests and every bounded stopping
time $\tau$, uniformly for $0\le h\le h_0$,
\[
 \mathbb E\left[
 \left|Z_{\tau+h}^\varepsilon-Z_\tau^\varepsilon\right|^4
 \mid\mathcal F_\tau\right]
 \le C\left(h^2+{h\over\mu_\varepsilon}\right)
\]

after stopping on one compact moment set; the stopped complement has
exponentially small probability.  Joint cumulants of order $k\ge3$ are
$O(\mu_\varepsilon^{1-k/2})$.
\end{lemma}

\begin{proof}
Condition on the microscopic state at $\tau$ and restart the deterministic
flow.  The B2 fixed-horizon trajectory expansion is uniform on compact
moment/trace sets and localizes every connected diagram to the interval
$[\tau,\tau+h]$.  Pairings of two contact insertions give $h^2$, while one
coincident contact diagonal gives the necessary $h/\mu_\varepsilon$ term.
Higher connected diagrams carry the usual power
$\mu^{1-k/2}$.  B2 exponential compact containment removes the stopped
complement.  This conditional estimate, rather than a deterministic-interval
bound alone, is the Aldous input.
\end{proof}

\begin{theorem}[Microscopic joint process convergence]
\label{thm:r13-b3-process}
The pair $(u^\varepsilon,\Xi^\varepsilon)$ converges in the weighted
trace-distribution Skorokhod space to the Gaussian pair of
Theorem~\ref{thm:r13-b3-joint}.  All density--density, density--contact, and
contact--contact covariances are the second derivatives of the limiting B2/B1
pressure.
\end{theorem}

\begin{proof}
Normal convergence of the marked pressure gives convergence of finite
cumulants.  Lemma~\ref{lem:r13-b3-cumulants} and Aldous--Mitoma give process
tightness, including the vanishing jump size.  Passing the exact microscopic
balance to the limit identifies the density equation, and the marked contact
coordinates identify the random measure $W$.  The initial Schur-complement
law comes from B1.  Uniqueness in
Theorem~\ref{thm:r13-b3-joint} identifies every subsequence.
\end{proof}

\subsection{The exact dual kernel and gauge}

Let $\mathcal Y$ be the weighted dual pair space with endpoint traces.  Define
\[
 \mathfrak D(p,\psi)=\Delta p+\psi,
 \qquad
 \mathfrak G r=(r,-\Delta r),
\]
and quotient additionally by endpoint/transport coboundaries whose complete
weak-balance pairing vanishes.  Let $C$ denote the initial constraint map.

\begin{theorem}[Exact covariance kernel]
\label{thm:r13-b3-kernel}
A dual class has zero covariance if and only if
\begin{enumerate}[label=(\roman*)]
\item its initial source lies in the span projected out by the B1 constraint
      Schur complement;
\item $\Delta p+\psi=0$ in $L^2(\mathfrak a)$; and
\item its remaining endpoint/transport pairing is a closed weak-balance
      coboundary.
\end{enumerate}
Thus the identifiable contact cotangent is the class of
$\Delta p+\psi$, and pure contact directions are not quotiented out.
\end{theorem}

\begin{proof}
The covariance is the sum of two nonnegative terms: the prepared initial
covariance and the $L^2(\mathfrak a)$ norm of
$\Delta p+\psi$ transported through the linear evolution.  Their vanishing
gives (i)--(ii).  The closed range of the weak balance operator on the Gelfand
triple from Theorem~\ref{thm:r13-b3-energy} identifies the residual annihilator
with (iii).  Conversely each listed direction pairs to zero with the initial
field and with $W$, so it has zero covariance.
\end{proof}

\subsection{The Cameron--Martin form and action tangent}

Let $\Sigma$ be the covariance operator of the quotient Gaussian process.
Define the closed extended form
\[
 \mathfrak q(z)=
 \begin{cases}
 \|\Sigma^{-1/2}z\|^2,
 &z\in\operatorname{Ran}\Sigma^{1/2},\\
 +\infty,&\text{otherwise}.
 \end{cases}
\]

\begin{theorem}[Second epi-derivative and Mosco limit]
\label{thm:r13-b3-mosco}
The second epi-derivative of the B2/B1 action at a regular exposed path is
$\tfrac12\mathfrak q$.  Finite-dimensional action Hessians Mosco-converge to
this closed form on the quotient Gelfand space.
\end{theorem}

\begin{proof}
For every finite projection, analytic pressure duality identifies the action
Hessian with the inverse of the projected covariance.  Realize all
projections as conditional expectations of the single Gaussian pair in
Theorem~\ref{thm:r13-b3-joint}; the corresponding Cameron--Martin forms are
increasing closed forms.  Monotone convergence of closed quadratic forms gives
the Mosco liminf and recovery on
$\operatorname{Ran}\Sigma^{1/2}$.  Density of finite projections in that range
and $+\infty$ outside it complete the proof.  The direct second variation of
the entropy action agrees on a core and hence has the same closure.
\end{proof}

\begin{theorem}[Round-thirteen B3 closure]
\label{thm:r13-b3-main}
The prepared hard-sphere tangent is a joint density--actual-contact Gaussian
process driven by one collision-space random measure.  Its exact dual noise
coefficient is $\Delta p+\psi$, its kernel is the closed balance/endpoint
gauge above, and its action tangent is the Cameron--Martin form on
$\operatorname{Ran}\Sigma^{1/2}$.
\end{theorem}

\begin{proof}
Combine Theorems~\ref{thm:r13-b3-joint},
\ref{thm:r13-b3-process}, \ref{thm:r13-b3-kernel}, and
\ref{thm:r13-b3-mosco}.
\end{proof}
